\section{Cup vanishing and the secondary obstruction}\label{sec:secondary}
The first equation $dp=-\alpha\cup\beta$ has an existence obstruction in $H^2(K,k)$. Once it vanishes, a further obstruction remains for the upper equation. The distinction is algebraic and can be proved without assuming an arithmetic realization.

Put $s(g)=\alpha(g)\beta(g)$. Expansion of the two cocycle laws gives
\begin{equation}\label{eq:ds}
ds=-(\alpha\cup\beta+\beta\cup\alpha).
\end{equation}
For a chosen primitive $p$, set $D_{12}=-p$ and $D_{23}=p-s$. They form a defining system with differentials $\alpha\cup\beta$ and $\beta\cup\alpha$. The associated twisted Massey type cocycle is
\begin{equation}\label{eq:mu}
\mu_p=\alpha\cup(p-s)-p\cup\alpha.
\end{equation}
Explicitly it equals
\[
\mu_p(g,h)=\alpha(g)(p(h)-\alpha(h)\beta(h))
-p(g)\chi(g)\alpha(h).
\]
Let $R_p$ denote the right side of \eqref{eq:q-equation}. Direct expansion of $f(g)=\alpha(g)p(g)$ gives
\begin{equation}\label{eq:secondaryidentity}
R_p+d_\chi f=\mu_p.
\end{equation}
Thus $[R_p]=[\mu_p]$ in $H^2(K,k(\chi))$. For this fixed $p$, the upper equation is soluble if and only if this class is zero.

Every other first primitive is $p+z$ with $z\in Z^1(K,k)$ an ordinary character. Formula \eqref{eq:mu} gives
\[
\mu_{p+z}-\mu_p=\alpha\cup z-z\cup\alpha.
\]
Graded commutativity in odd characteristic makes its cohomology class $2[\alpha]\cup[z]$. Therefore the choice independent object at this stage is the affine coset
\begin{equation}\label{eq:masseycoset}
[\mu_p]+2[\alpha]\cup H^1(K,k).
\end{equation}
There exists some first primitive for which the upper equation is soluble precisely when zero belongs to this coset. There is not one canonical class independent of the primitive without taking this quotient by indeterminacy. With restricted ramification or prescribed Stage P local conditions, only the corresponding allowed character subspace may vary $p$. After the final Stage P reduction this is $\ker A$, so the larger absolute indeterminacy cannot discharge those affine conditions.

\subsection{An exact control separating the two stages}
Take the finite group $U_3(\F_7)$ with coordinates $(a,b,c)$ and multiplication
\[
(a,b,c)(A,B,C)=(a+A,b+B,c+C+aB).
\]
Use trivial cyclotomic action, $\alpha=a$, $\beta=b$, and $p=c$. Then $dp=-\alpha\cup\beta$, so the first obstruction vanishes. The elements $x=(1,0,0)$ and $z=(0,0,1)$ commute, but
\[
R_p(x,z)=2,\qquad R_p(z,x)=0.
\]
For every ordinary cochain $q$, commuting arguments give $dq(x,z)=dq(z,x)$. Hence no $q$ solves the second equation. This is a finite group control, not an absolute Galois group or a Frey representation. It proves that cup zero alone does not formally imply upper solvability.

\subsection{The fixed primitive local to global theorem}
For $K=\Q(\sqrt5)$, localization on $H^2(K,k(\chi))$ is injective. To see the coefficient translation, restrict to $L=K(\zeta_7)$, whose degree over $K$ is six. Restriction followed by corestriction multiplies by six, which is invertible on seven primary cohomology. Over $L$, the module is two trivial $\F_7$ copies and $\mu_7$ is present. Kummer theory identifies each copy of $H^2(L,\F_7)$ with $\operatorname{Br}(L)[7]$. Brauer--Hasse--Noether gives local to global injectivity for the Brauer group. A class zero at every completion over $K$ is therefore zero over $L$ and then over $K$ \cite{milneCFT2020}. The exact leaves are II.1.22 and II.1.30 for Kummer and restriction with corestriction, and VIII.4.2 for the Brauer sequence.

At a fixed place where beta is the coboundary of $\lambda_v$ and $p=p_v^{\mathrm{fix}}$, the explicit local upper primitive
\[
q_v^{\mathrm{fix}}(g)=-\lambda_v\chi(g)^{-1}\alpha(g)^2
\]
satisfies $d_\chi q_v^{\mathrm{fix}}=R_p$ by direct expansion. At support, $\alpha=0$ also kills the local upper right side. These facts prevent treating a formal secondary obstruction as an established nonzero obstruction in the actual local branches. Full prescribed upper values and geometric compatibility remain separate.

The local to global theorem uses the same global $p$ at every place. Independently chosen local primitives can have vanishing secondary classes without supplying one such global primitive. Even once $q$ exists globally, its local values must match the actual toric, endpoint, semilinear, and fixed seven conditions. Those conditions form another affine localization problem.

General Massey vanishing theorems do not erase these qualifications. The all prime theorem of Min\'{a}\v{c} and T\^{a}n concerns ordinary products over absolute Galois groups with trivial $\F_p$ coefficients \cite{minac-tan2016}. The JEMS theorem cited in some older discussion has a mod two statement \cite{minac-tan2017}. The 2026 strong vanishing theorem likewise concerns its stated ordinary coefficient setting and may introduce additional ramification \cite{maire-et-al2026}. None is a theorem that the one twisted defining system prescribed here, on a fixed ramification set and with actual affine local targets, has a compatible integral lift.
